\chapter{The Filesystem Hierarchy, Inode Architecture \& File Management}
\label{chap:filesystem-management}

\section{The Filesystem Hierarchy Standard (FHS)}
In contrast to operating systems that employ multiple drive letters (\texttt{C:}, \texttt{D:}), Linux arranges all storage devices, partitions, pseudo-filesystems, and network mounts beneath a single inverted tree rooted at the forward slash (\texttt{/}). The layout is codified by the \textbf{Filesystem Hierarchy Standard (FHS)}.

\begin{figure}[htbp]
\centering
\begin{tikzpicture}[node distance=1.0cm,
    dir/.style={rectangle, draw=brandnavy, fill=boxnavybg, rounded corners=2pt, font=\small\ttfamily\bfseries, minimum width=1.8cm, minimum height=0.7cm},
    edge/.style={thick, draw=brandnavy}]
    \node [dir] (root) {/ (Root)};
    \node [dir, below left=1.2cm and 3.5cm of root] (bin) {/bin};
    \node [dir, below left=1.2cm and 2.0cm of root] (etc) {/etc};
    \node [dir, below left=1.2cm and 0.5cm of root] (home) {/home};
    \node [dir, below right=1.2cm and 0.5cm of root] (dev) {/dev};
    \node [dir, below right=1.2cm and 2.0cm of root] (proc) {/proc};
    \node [dir, below right=1.2cm and 3.5cm of root] (var) {/var};

    \draw [edge] (root) -- (bin);
    \draw [edge] (root) -- (etc);
    \draw [edge] (root) -- (home);
    \draw [edge] (root) -- (dev);
    \draw [edge] (root) -- (proc);
    \draw [edge] (root) -- (var);
\end{tikzpicture}
\caption{The Root Branching Structure under the Filesystem Hierarchy Standard.}
\label{fig:fhs-tree}
\end{figure}

\begin{table}[htbp]
\centering
\small
\begin{tabularx}{\linewidth}{l X}
\toprule
\textbf{Directory Path} & \textbf{FHS Purpose \& Structural Contents} \\
\midrule
\texttt{/} & The Root Directory; the supreme parent of the entire hierarchy. \\
\texttt{/bin} & Essential user command binaries required in single-user recovery mode (\texttt{cat}, \texttt{ls}, \texttt{cp}). \\
\texttt{/sbin} & Essential system binaries for superuser administration (\texttt{fdisk}, \texttt{init}, \texttt{ip}). \\
\texttt{/etc} & Host-specific system-wide configuration files and startup scripts (\texttt{passwd}, \texttt{fstab}). \\
\texttt{/home} & User home directories housing personal configurations and documents (\texttt{/home/asha}). \\
\texttt{/root} & Dedicated home directory of the superuser (root). \\
\texttt{/dev} & Device special files representing physical/virtual hardware (\texttt{/dev/sda}, \texttt{/dev/null}, \texttt{/dev/urandom}). \\
\texttt{/proc} & Virtual pseudo-filesystem exposing real-time kernel data structures and process state. \\
\texttt{/sys} & Unified sysfs pseudo-filesystem exposing hardware buses and kernel device drivers. \\
\texttt{/var} & Variable data files: system logs (\texttt{/var/log}), spool queues, and mail stores. \\
\texttt{/tmp} & Temporary storage files accessible by all users; automatically emptied on reboot. \\
\texttt{/usr} & User secondary hierarchy containing sharable, read-only multi-user software (\texttt{/usr/bin}). \\
\texttt{/opt} & Add-on commercial or third-party proprietary application software packages. \\
\texttt{/boot} & Kernel images (\texttt{vmlinuz}), initramfs ramdisks, and GRUB bootloader configuration. \\
\texttt{/lib} & Shared dynamic library dependencies (\texttt{.so} files) and kernel modules. \\
\bottomrule
\end{tabularx}
\caption{Linux Filesystem Hierarchy Standard (FHS) Directory Directory Specifications.}
\label{tab:fhs}
\end{table}

\section{Path Resolution: Absolute vs Relative Paths}
Every file in Linux is uniquely addressed through a path string:
\begin{itemize}[leftmargin=1.5em]
    \item \textbf{Absolute Path:} Defined from the root directory (\texttt{/}). It invariably commences with a forward slash (e.g., \texttt{/home/student/projects/main.c}). It resolves unambiguously regardless of the working directory.
    \item \textbf{Relative Path:} Defined with respect to the current working directory. It never begins with a leading forward slash:
    \begin{itemize}
        \item \texttt{.} refers explicitly to the current working directory.
        \item \texttt{..} refers explicitly to the immediate parent directory.
        \item \texttt{\textasciitilde} resolves to the active user's home directory.
        \item \texttt{-} represents the previous working directory (\texttt{OLDPWD}).
    \end{itemize}
\end{itemize}

\section{The Inode Architecture \& Hard vs Symbolic Links}
In Linux filesystems (such as \texttt{ext4}), a file is strictly split into two entities: \textbf{Directory Entry} (File name and inode pointer) and \textbf{Inode} (Index Node containing metadata and pointers to data disk blocks).

\begin{definition}[Inode Metadata Attributes]
The inode data structure records:
\begin{enumerate}[leftmargin=1.5em]
    \item File type (regular file, directory, symlink, block device, FIFO).
    \item Permissions and access control bits (mode flags).
    \item Owner User ID (UID) and Owner Group ID (GID).
    \item File byte size.
    \item Hard link reference count.
    \item Timestamps: Access time (\texttt{atime}), Modification time (\texttt{mtime}), and Inode Change time (\texttt{ctime}).
    \item Direct and indirect block pointers addressing physical storage sectors.
\end{enumerate}
\textit{Crucial Principle:} An inode does \textbf{not} contain the file name. The file name is stored exclusively inside the parent directory table alongside its associated inode number.
\end{definition}

\subsection{Hard Links versus Symbolic (Soft) Links}
\begin{table}[htbp]
\centering
\small
\begin{tabularx}{\linewidth}{l X X}
\toprule
\textbf{Feature Dimension} & \textbf{Hard Link (\texttt{ln target link})} & \textbf{Symbolic Link (\texttt{ln -s target link})} \\
\midrule
\textbf{Inode Allocation} & Shares the \textit{exact same} inode number as the original target file. & Allocates a \textit{new, distinct} inode with file type \texttt{l}. \\
\textbf{Data Content} & Directly points to the target file's physical data blocks. & Contains a text string storing the pathname of the target. \\
\textbf{Target Deletion} & Data remains intact; link count decrements by 1; data freed only when link count reaches 0. & Becomes a broken (dangling) link pointing to a non-existent target. \\
\textbf{Cross-Filesystem} & \textbf{Cannot} span across distinct partitions or filesystems (inode IDs are local). & \textbf{Can} cross distinct storage drives, filesystems, and network shares. \\
\textbf{Directory Linking} & Prohibited for ordinary users to prevent recursive cyclic infinite directory loops. & Permitted and widely utilized to reference directories. \\
\bottomrule
\end{tabularx}
\caption{Comparison Between Hard Links and Symbolic Links.}
\label{tab:links}
\end{table}

\section{Core Directory \& File Operations Primitives}
\begin{itemize}[leftmargin=1.5em]
    \item \texttt{pwd}: Print Working Directory. Outputs the absolute path of the current directory.
    \item \texttt{cd [dir]}: Change Directory. \texttt{cd} with no arguments returns to \texttt{\$HOME}; \texttt{cd -} toggles to the previous directory.
    \item \texttt{ls [options] [path]}: Lists directory contents.
    \begin{itemize}
        \item \texttt{-l}: Long listing format (permissions, links, owner, group, size, mtime, name).
        \item \texttt{-a}: Show all files, including hidden files (commencing with a dot \texttt{.}).
        \item \texttt{-h}: Human-readable sizes (K, M, G).
        \item \texttt{-t}: Sort by modification time (most recent first).
        \item \texttt{-r}: Reverse the sort order.
        \item \texttt{-i}: Display inode index numbers.
    \end{itemize}
    \item \texttt{mkdir -p a/b/c}: Creates nested directory hierarchies without error.
    \item \texttt{rmdir}: Removes empty directories only.
    \item \texttt{touch}: Creates empty files if non-existent, or updates \texttt{atime} and \texttt{mtime} timestamps.
    \item \texttt{cp -r -p src dest}: Recursively copies directories (\texttt{-r}) preserving timestamps and permissions (\texttt{-p}).
    \item \texttt{mv}: Moves files or renames filenames atomically within the same filesystem.
    \item \texttt{rm -rf}: Recursively and forcefully deletes directories and files without interactive prompting.
\end{itemize}

\section{Deep Search: The find Engine}
Unlike \texttt{locate} which relies on a pre-indexed database, \textbf{find} navigates physical directory inodes in real-time.
\begin{commandbox}{find --- Search Files in a Directory Hierarchy}
\textbf{Syntax:} \texttt{find [starting-point] [criteria] [action]}
\begin{itemize}[leftmargin=1.5em]
    \item \texttt{-name "pattern"}: Case-sensitive name matching with shell wildcards.
    \item \texttt{-iname "pattern"}: Case-insensitive name matching.
    \item \texttt{-type f|d|l}: Matches regular files (\texttt{f}), directories (\texttt{d}), or symlinks (\texttt{l}).
    \item \texttt{-size +100M}: Matches files exceeding 100 Megabytes in size.
    \item \texttt{-mtime -7}: Modified within the past 7 days.
    \item \texttt{-perm 644}: Exact octal permission matching.
    \item \texttt{-exec cmd \{\} \textbackslash;}: Executes \texttt{cmd} on each matching file sequentially.
    \item \texttt{-exec cmd \{\} +}: Batches matching filenames together into a single command invocation (high performance).
\end{itemize}
\end{commandbox}

\begin{terminalbox}[Advanced find Practical Invocations]
$ find /var/log -type f -name "*.log" -size +10M
/var/log/syslog.1
/var/log/nginx/access.log

$ find /home/user/docs -type f -mtime -3 -exec cp {} /backup/recent/ +
\end{terminalbox}

\section{Directory Visualization \& Disk Space Analysis}
\begin{itemize}[leftmargin=1.5em]
    \item \textbf{tree:} Produces a visual tree representation. \texttt{tree -L 2 -d} limits recursion depth to 2 levels and restricts display to directories only.
    \item \textbf{history:} Inspects the interactive command history list. \texttt{!!} repeats the preceding command; \texttt{!n} executes command number $n$; \texttt{!-2} executes the second-to-last command; \texttt{!str} executes the most recent command beginning with \texttt{str}.
    \item \textbf{df -h:} Reports filesystem disk space usage, total capacity, and mount points in human-readable units. The \texttt{-i} option audits inode capacity.
    \item \textbf{du -sh *:} Estimates disk space usage of files and directories within the current working folder.
\end{itemize}

\section{Chapter Review \& Analytical Questions}
\begin{enumerate}[leftmargin=1.5em]
    \item State the architectural purpose of \texttt{/proc} and explain why files within it consume 0 bytes of physical disk space.
    \item What is the functional impact on a file when its hard link count drops from 2 to 1? What occurs when it drops from 1 to 0?
    \item Why cannot a hard link be established across two distinct physical disk partitions?
    \item Construct a \texttt{find} command that locates all regular files ending with \texttt{.tmp} inside \texttt{/tmp} older than 30 days and removes them safely.
    \item Differentiate between \texttt{df -h} and \texttt{du -sh}. In what situation can \texttt{df} report a disk as 100\% full even when \texttt{du} reports ample free storage?
\end{enumerate}
